\documentclass[10pt,twocolumn]{article}
\usepackage[margin=0.75in]{geometry}
\usepackage{times}
\usepackage{amsmath,amssymb}
\usepackage{booktabs}
\usepackage{graphicx}
\usepackage{hyperref}
\usepackage{xcolor}
\usepackage{multirow}
\usepackage{array}
\usepackage{enumitem}
\usepackage{microtype}

\hypersetup{
  colorlinks=true,
  linkcolor=blue!70!black,
  citecolor=green!50!black,
  urlcolor=blue!70!black
}

\newcommand{\dwb}{\textsc{DWB}}
\newcommand{\tq}{\textsc{TurboQuant}}
\newcommand{\dwbtq}{\textsc{DWB-TurboQuant}}
\newcommand{\eg}{\textit{e.g.,}\xspace}
\newcommand{\ie}{\textit{i.e.,}\xspace}
\newcommand{\etal}{\textit{et al.}\xspace}

\title{\textbf{Independent Verification of \textit{Don't Waste Bits!}:}\\
Scale-Dependent INT4 Losslessness and a DWB-TurboQuant Extension}

\author{
  themoddedcube / LonghornSilicon\\
  \texttt{github.com/LonghornSilicon/dont-waste-bits}
}

\date{April 2026}

\begin{document}
\maketitle

\begin{abstract}
We present an independent reproduction of ``Don't Waste Bits!''
\cite{haeri2026dwb} (CVPR 2026), which proposes adaptive per-token
KV-cache quantization for on-device LLMs via a learned controller that
assigns \{2,4,8,16\}-bit precision to each token.
Since the original code is not yet public (expected June 2026), we
re-implement the method from paper equations and identify seven
methodological insights.
Our key finding is that symmetric per-tensor INT4 KV quantization
exhibits \textbf{scale-dependent losslessness}: at 135M and 360M
parameters (15 attention heads), INT4 is nearly lossless
($\approx$0pp gap vs FP16), but at 1.7B parameters (32 heads) it shows
genuine $\approx$10pp degradation matching the paper's stated baseline.
The paper's 33.6\% static INT4 baseline for smaller models is reproduced
only with a non-standard 8-level quantization (scale=max/3, equivalent
to INT3 precision stored in 4-bit), confirmed causally by ablation.
We mechanistically verify the losslessness via
$\epsilon_\text{eff} = \epsilon_\text{rel} \times r_\text{cancel}$:
8.1\% at 360M (below threshold $\Rightarrow$ lossless) vs.\
12.4\% at 1.7B (above threshold $\Rightarrow$ 10pp loss).
For the DWB adaptive controller (H3), we find an implementation gap:
our v1 achieves 33.8\% vs.\ the paper's 41.2\% ($-$7.4pp, outside
CI$\pm$4.4pp), while our v2 (better-trained) achieves 37.0\% but at
1.68$\times$ more bits, revealing a dual-objective tension—quartile
classification training cannot simultaneously optimize accuracy and
compression.
As a novel contribution, we propose \dwbtq, which routes
DWB's 2-bit tokens through PolarQuant (per-head Walsh-Hadamard
rotation) rather than naive scalar INT2, recovering
$+$2pp on HellaSwag and $+$3pp on ARC-Challenge at identical
compression (5.05 avg bits).
\end{abstract}

%-----------------------------------------------------------------------
\section{Introduction}
\label{sec:intro}
%-----------------------------------------------------------------------

Large language models deployed on-device face a fundamental memory
bandwidth bottleneck during autoregressive decoding: the KV cache grows
linearly with sequence length and must be read from memory on every
step \cite{pope2023efficiently}.
Quantizing the KV cache to lower bit-widths reduces this bottleneck at
the cost of accuracy \cite{liu2024kivi,hooper2024kvquant}.

``Don't Waste Bits!'' (DWB) \cite{haeri2026dwb}, accepted at CVPR 2026,
proposes a per-token adaptive approach: a lightweight 3-layer MLP
controller predicts a precision assignment $b_t \in \{2,4,8,16\}$ for
each token using four signals—entropy $H_t$, rarity $R_t$, attention
variance $V_t$, and prediction confidence $C_t$.
The paper reports that for SmolLM-360M \cite{allal2024smollm} on HellaSwag
\cite{zellers2019hellaswag}: (H1) 17.75\% latency reduction over static
4-bit KV; (H2) +7.60pp accuracy over static 4-bit KV; (H3) within
0.30pp of FP16 baseline.
The original code is not yet public, making independent verification both
necessary and valuable.

\paragraph{Contributions.}
\begin{enumerate}[leftmargin=*,topsep=2pt,itemsep=0pt]
  \item Re-implementation of DWB from paper equations with seven
        methodological discoveries (\S\ref{sec:insights}).
  \item Scale-dependent INT4 losslessness: mechanistic explanation
        via effective residual error across SmolLM-135M/360M/1.7B
        (\S\ref{sec:results}).
  \item \dwbtq: routing DWB's 2-bit tier through PolarQuant
        \cite{turboQuant2026}, recovering $+$2--3pp at identical
        compression (\S\ref{sec:turboquant}).
\end{enumerate}

%-----------------------------------------------------------------------
\section{Background}
\label{sec:background}
%-----------------------------------------------------------------------

\paragraph{KV cache quantization.}
KV cache quantization stores keys and values at reduced precision
\cite{liu2024kivi,hooper2024kvquant,yuan2024kvcache}.
Standard INT4 (16 quantization levels, scale = max/7) minimally impacts
attention scores in practice \cite{liu2024kivi}.
DWB \cite{haeri2026dwb} improves over static quantization by
importance-aware per-token assignment.

\paragraph{SmolLM model family.}
SmolLM \cite{allal2024smollm} is a family of small causal LMs:
135M, 360M, and 1.7B parameters.
SmolLM-360M uses 32 transformer layers, 15 attention heads, and hidden
dimension 960; SmolLM-1.7B uses 24 layers, 32 heads, and hidden
dimension 2048.

\paragraph{HellaSwag evaluation.}
HellaSwag \cite{zellers2019hellaswag} is a commonsense NLI benchmark
with 4-way completion.
The paper uses \emph{unnormalized} log-likelihood accuracy (\texttt{acc},
$\approx$42\% for SmolLM-360M at FP16), not length-normalized
(\texttt{acc\_norm}, $\approx$54\%).
This distinction is critical for reproducing reported numbers
(\S\ref{sec:insights}).

\paragraph{TurboQuant / PolarQuant.}
TurboQuant \cite{turboQuant2026} (ICLR 2026) applies a per-head
Walsh-Hadamard Transform (WHT) rotation to KV tensors before uniform
scalar quantization, decorrelating activations and reducing effective
quantization error.
The rotation is orthogonal and exactly invertible, requiring no extra
bits.

%-----------------------------------------------------------------------
\section{Re-Implementation}
\label{sec:method}
%-----------------------------------------------------------------------

\paragraph{Controller architecture.}
We re-implement the DWB controller as:
\[
  f_\theta: \mathbb{R}^4 \to \mathbb{R}^4, \quad
  f_\theta = \text{Linear}(128,4) \circ \text{ReLU} \circ
  \text{Linear}(128,128) \circ \text{ReLU} \circ \text{Linear}(4,128),
\]
totalling 33,540 parameters.
For each token $t$, the controller takes signals
$[H_t, R_t, V_t, C_t]$ (Eqs.\ 14--17 of \cite{haeri2026dwb}) and
outputs logits over $\{2,4,8,16\}$-bit classes.

\paragraph{Signals.}
$H_t = -\sum_v p_v \log p_v$ (next-token entropy);
$R_t$ = inverse token frequency;
$V_t$ = variance of attention weights across heads (requires
\texttt{output\_attentions=True});
$C_t$ = max softmax probability.

\paragraph{Training loss.}
We follow Eq.\ 28 of \cite{haeri2026dwb}:
$\mathcal{L} = \alpha \cdot \mathcal{L}_\text{CE} +
\beta \cdot \mathcal{L}_\text{latency} +
\gamma \cdot \mathcal{L}_\text{quality}$,
with $\alpha=1, \beta=0.1, \gamma=0.1$.
The latency term penalizes high-bit assignments with costs
$[0.5, 1.0, 2.0, 4.0]$ for $\{2,4,8,16\}$-bit respectively.

\paragraph{KV cache hooks.}
\textbf{Critical finding (Insight 2):} transformers $\geq$5.0 uses
\texttt{DynamicCache} objects; hooks on attention output tuples silently
fail.
Correct approach: register forward hooks on \texttt{k\_proj} and
\texttt{v\_proj} Linear submodules directly.
For SmolLM-360M this places 64 hooks (32 layers $\times$ 2).

%-----------------------------------------------------------------------
\section{Evaluation Setup}
\label{sec:setup}
%-----------------------------------------------------------------------

We evaluate on HellaSwag validation using unnormalized log-likelihood
accuracy (\texttt{acc}).
Primary model: SmolLM-360M \cite{allal2024smollm}.
Cross-model validation on SmolLM-135M and SmolLM-1.7B.
Hardware: CPU (accuracy experiments); latency experiments (H1) require
RTX 4090 (pending).
All results use float32.

%-----------------------------------------------------------------------
\section{Results}
\label{sec:results}
%-----------------------------------------------------------------------

\subsection{Baseline Verification}

Table~\ref{tab:main} shows our SmolLM-360M results vs.\ paper Table~3.

\begin{table}[h]
\centering
\caption{SmolLM-360M on HellaSwag (\texttt{acc}).
CI $\pm$4.4pp at $n$=500.}
\label{tab:main}
\footnotesize
\begin{tabular}{lccc}
\toprule
Condition & $n$ & Ours & Paper \\
\midrule
FP16 & 500 & \textbf{42.6\%} & 41.5\% \\
KV-4bit per-tensor & 500 & 41.6\% & 33.6\% \\
KV-4bit per-token & 500 & 41.2\% & 33.6\% \\
KV-4bit asymmetric & 200 & 42.5\% & 33.6\% \\
\textbf{int4\_int3range (8-level)} & 100 & \textbf{33.0\%} & \textbf{33.6\%} \\
KV-2bit & 200 & 25.0\% & --- \\
DWB v1 (val\_acc=0.366) & 500 & 33.8\% & 41.2\% \\
DWB v2 (val\_acc=0.446) & 500 & 37.0\% & 41.2\% \\
\bottomrule
\end{tabular}
\end{table}

The FP16 baseline (42.6\%) matches the paper's 41.5\% within noise ($+$1.1pp).
All standard INT4 variants are near-lossless at 360M—a striking result
that contradicts the paper's reported 33.6\% baseline.
The 33.6\% baseline is reproduced only by \texttt{int4\_int3range}:
scale=max/3 with 8 effective levels, equivalent to INT3 precision in
4-bit storage.

\subsection{Scale-Dependent INT4 Losslessness}

\begin{table}[h]
\centering
\caption{H4 cross-model validation.}
\label{tab:h4}
\footnotesize
\begin{tabular}{llcc}
\toprule
Model & Condition & Ours & Paper \\
\midrule
\multirow{3}{*}{135M} & FP16 & 40.0\% & 37.2\% \\
                       & Std INT4 & 39.0\% & 33.6\% \\
                       & int4\_int3range & \textbf{32.0\%} & \textbf{33.6\%} \\
\midrule
\multirow{3}{*}{1.7B} & FP16 & 50.0\% & 49.0\% \\
                       & \textbf{Std INT4} & \textbf{40.0\%} & \textbf{41.1\%} \\
                       & int4\_int3range & 32.0\% & 41.1\% \\
\bottomrule
\end{tabular}
\end{table}

A critical reversal appears at 1.7B: standard INT4 now \emph{matches} the
paper's baseline (40.0\% vs.\ 41.1\%, $-$1.1pp), while int4\_int3range
over-degrades ($-$9.1pp).
At 135M/360M (15 heads), standard INT4 is lossless ($\leq$1.4pp gap vs.\
FP16).
At 1.7B (32 heads), it degrades $\approx$10pp.

\paragraph{Mechanistic explanation.}
We directly measure the quantization error in K/V tensors (20 examples each):

\begin{table}[h]
\centering
\caption{INT4 error analysis across scales.}
\label{tab:mech}
\footnotesize
\begin{tabular}{lcc}
\toprule
Metric & 360M & 1.7B \\
\midrule
Symmetry (mean/std) & 0.0027 & 0.0006 \\
Relative error $\epsilon_\text{rel}$ & 26.95\% & 35.31\% \\
Cancellation ratio $r_\text{cancel}$ & 0.30 & 0.35 \\
\textbf{Eff.\ residual} $\epsilon_\text{eff}$ & \textbf{8.1\%} & \textbf{12.4\%} \\
Accuracy impact & $\approx$0pp & $\approx$10pp \\
\bottomrule
\end{tabular}
\end{table}

Both scales have near-zero-mean errors (symmetry $\approx$0),
confirming the same cancellation mechanism operates at both scales.
The effective residual $\epsilon_\text{eff} = \epsilon_\text{rel} \times
r_\text{cancel}$ captures what survives after cancellation in the
attention weighted sum:
\begin{align*}
  \text{360M}: \quad & 26.95\% \times 0.30 = 8.1\% \quad \Rightarrow \text{lossless} \\
  \text{1.7B}: \quad & 35.31\% \times 0.35 = 12.4\% \quad \Rightarrow -10\text{pp loss}
\end{align*}
The decision threshold lies between 8.1\% and 12.4\%.
\textbf{Root cause}: SmolLM-1.7B's larger hidden dimension (2048 vs.\ 960)
produces higher-variance KV tensors; at the same scale divisor (max/7),
this yields larger relative quantization errors.

\paragraph{Implication for H2.}
The paper's H2 claim (+7.6pp over static INT4) is most strongly supported
at 1.7B, where standard INT4 \emph{genuinely} degrades.
At 135M/360M, the +7.6pp improvement is over a sub-standard 8-level
baseline—standard INT4 at these scales already matches FP16.

\subsection{INT4 Baseline Ablation}

We isolate the effect of step size vs.\ range clipping via a 5-condition
ablation (50 samples):

\begin{table}[h]
\centering
\caption{INT4 ablation: step size vs.\ range clipping.}
\label{tab:ablation}
\footnotesize
\begin{tabular}{lccc}
\toprule
Condition & Scale & Clamp & Acc \\
\midrule
A: Standard & max/7 & $(-8,7)$ & 46.0\% \\
E: Intermediate & max/5 & $(-8,7)$ & 42.0\% \\
D: Range only & max/7 & $(-4,3)$ & 38.0\% \\
C: Coarse step & max/3 & $(-8,7)$ & 28.0\% \\
B: int4\_int3range & max/3 & $(-4,3)$ & 28.0\% \\
\bottomrule
\end{tabular}
\end{table}

B $=$ C (both 28\%): range clipping contributes \textbf{0pp} once the
step is coarse.
All $-$18pp degradation in the paper's baseline is caused by the coarse
scale (max/3), not range clipping.
Standard max/7 (16 levels) is safely within the lossless regime.

\subsection{DWB Controller Evaluation (H3)}

\begin{table}[h]
\centering
\caption{DWB controller sensitivity study.}
\label{tab:dwb}
\footnotesize
\begin{tabular}{lccccc}
\toprule
Controller & Train & Ep. & val\_acc & avg\_bits & DWB acc \\
\midrule
v1 & 100 & 5 & 0.366 & 5.03 & 33.8\% \\
v2 & 500 & 10 & 0.446 & 8.47 & 37.0\% \\
\textit{Paper} & --- & --- & --- & \textit{5.05} & \textit{41.2\%} \\
\bottomrule
\end{tabular}
\end{table}

v1 (33.8\% at 5.03 avg\_bits) performs \emph{worse} than standard INT4
(41.6\% at 4.0 bits)—the controller mislabels important tokens as 2-bit.
v2 (better-trained) improves to 37.0\% but requires 8.47 avg\_bits (vs.\
paper's 5.05)—a 68\% increase.
The paper achieves \emph{both} targets simultaneously via its compound
loss with latency penalty; our quartile-classification approach
optimizes classification accuracy alone.
The v2 bit distribution is bimodal: $\{2: 41.7\%, 4: 9.3\%, 8: 7.2\%, 16: 41.8\%\}$—the controller
over-polarizes between unimportant (2-bit) and critical (16-bit) tiers,
losing the intermediate compression range.

\subsection{H3 Beta-Sweep: Can Latency Penalty Resolve the Tension?}
\label{sec:beta_sweep}

We test whether the dual-objective tension is simply a hyperparameter
issue—i.e., whether increasing the latency penalty weight $\beta$ in
$\mathcal{L} = \alpha\mathcal{L}_\text{CE} + \beta\mathcal{L}_\text{latency} + \gamma\mathcal{L}_\text{quality}$
can drive avg\_bits toward the paper's 5.05 without collapsing accuracy.
We sweep $\beta \in \{0.1, 0.5, 1.0, 2.0\}$ with 100 train/eval samples
and 5 epochs, holding $\alpha=1, \gamma=0.1$ fixed.

\begin{table}[h]
\centering
\caption{$\beta$ sweep: latency penalty weight vs.\ accuracy and
compression. No $\beta$ achieves the paper's 41.2\% at 5.05 bits.
At $\beta \geq 1.0$ the controller assigns every token to 2-bit
(100\% 2-bit), collapsing accuracy to 32\% and val\_acc below random
(0.25 for 4 classes). Full 4-beta sweep (100 train/eval, 5 epochs).}
\label{tab:beta}
\footnotesize
\begin{tabular}{ccccc}
\toprule
$\beta$ & Acc & avg\_bits & val\_acc & Bit dist (\% 2b/4b/8b/16b) \\
\midrule
0.1 & 39.0\% & 5.30 & 0.407 & 38.5/37.6/9.9/14.0 \\
0.5 & 39.0\% & 3.92 & 0.341 & 37.0/57.4/0.0/5.5 \\
1.0 & 32.0\% & \textbf{2.00} & 0.257 & 100/0/0/0 \\
2.0 & 32.0\% & \textbf{2.00} & 0.227 & 100/0/0/0 \\
\midrule
\textit{Paper} & \textit{41.2\%} & \textit{5.05} & --- & undisclosed \\
\bottomrule
\end{tabular}
\end{table}

\textbf{Finding}: The dual-objective tension is \emph{not} a
hyperparameter issue.
At $\beta=0.1$, avg\_bits = 5.30 $\approx$ paper's 5.05, yet accuracy
(39.0\%) remains 2.2pp below the paper's 41.2\%.
At $\beta=0.5$, avg\_bits drops to 3.92 while accuracy is unchanged
(39.0\%)---higher latency penalty reduces bits with no accuracy benefit.
At $\beta \geq 1.0$, the controller collapses entirely: every token is
assigned 2-bit (avg\_bits = 2.00, bit\_dist = \{2: 100\%\}), accuracy
drops to 32.0\%, and val\_acc (0.257 at $\beta=1.0$; 0.227 at $\beta=2.0$)
falls below the random baseline of 0.25---the latency penalty completely
overwhelms the classification objective.
There is no $\beta$ that simultaneously achieves the paper's 41.2\%
accuracy at 5.05 avg\_bits; end-to-end compound-loss training is
qualitatively different from our quartile-classification approach.
We will update this table when results are available.

%-----------------------------------------------------------------------
\section{Methodological Insights}
\label{sec:insights}
%-----------------------------------------------------------------------

\paragraph{Insight 1: Evaluation metric.}
\texttt{acc\_norm} (length-normalized, lm-eval default) gives $\approx$54\%
for SmolLM-360M—matching SmolLM-1.7B's FP16 score rather than the
paper's 41.5\%.
The paper uses unnormalized \texttt{acc} ($\approx$42\%).
This single discrepancy would make every comparison meaningless.

\paragraph{Insight 2: KV cache hooks (transformers 5.x).}
transformers $\geq$5.0 wraps KV tensors in a \texttt{DynamicCache} object.
Hooks registered on attention module outputs silently intercept an empty
or wrong tensor.
\textbf{Fix}: hook \texttt{k\_proj} and \texttt{v\_proj} Linear submodules
directly (32 layers $\times$ 2 = 64 hooks for SmolLM-360M).
Validation: KV-2bit reduces accuracy to 25\% (near-chance), confirming
hooks fire correctly.

\paragraph{Insight 3: sdpa blocks \texttt{output\_attentions}.}
transformers 5.x uses scaled dot-product attention (sdpa) by default,
which does not support \texttt{output\_attentions=True} and silently
returns an empty tuple.
\textbf{Fix}: reload with \texttt{attn\_implementation=`eager'} for DWB
signal extraction pass.

\paragraph{Insight 4: Scale-dependent INT4 losslessness.}
See \S\ref{sec:results}.
The losslessness mechanism is: symmetric (zero-mean) quantization errors
partially cancel in the attention weighted sum; the fraction that
survives ($\epsilon_\text{eff}$) determines whether accuracy is
preserved.
This mechanism holds at all scales but crosses a critical threshold
between 360M and 1.7B.

\paragraph{Insight 5: Paper's baseline uses 8-level INT4.}
\texttt{int4\_int3range} (scale=max/3, 8 levels) reproduces the 33.6\%
baseline at 135M/360M ($-$0.6pp).
Standard INT4 (scale=max/7, 16 levels) is lossless.
The paper's ``Static 4-bit KV'' is INT3 precision in 4-bit storage.

\paragraph{Insight 6: Controller behavior—C$_t$ dominates.}
Analysis on 1,484 tokens (50 examples).
Cohen's $d$ between 2-bit and 16-bit tiers:
$C_t = 4.55$ (primary driver),
$H_t = 4.09$ (strong secondary),
$V_t = 1.42$ (moderate),
$R_t = 0.52$ (near-uninformative on HellaSwag vocabulary).
The controller approximates: \emph{high-confidence, low-entropy tokens
get more bits.}
$R_t$ (inverse frequency) barely discriminates despite its inclusion in
the paper's training loss.

\paragraph{Insight 7: Dual-objective tension.}
Our quartile-classification training optimizes controller classification
accuracy but not compression.
A better-trained controller (v2) becomes conservative—assigning 16-bit
to ambiguous tokens and 2-bit to clear ones, producing a bimodal
distribution.
The paper's compound loss $\mathcal{L} = \alpha\mathcal{L}_\text{CE} +
\beta\mathcal{L}_\text{latency} + \gamma\mathcal{L}_\text{quality}$
explicitly penalizes high average bit-width, forcing the controller to
find the accuracy/compression Pareto front.
Reproducing H3 fully requires the paper's undisclosed training details.

%-----------------------------------------------------------------------
\section{DWB-TurboQuant Extension}
\label{sec:turboquant}
%-----------------------------------------------------------------------

\paragraph{Motivation.}
DWB assigns 42--57\% of tokens to 2-bit scalar quantization.
Scalar INT2 is highly destructive (25.0\% accuracy alone).
PolarQuant (from TurboQuant \cite{turboQuant2026}) applies a per-head
Walsh-Hadamard Transform before scalar quantization, decorrelating
activations and reducing effective error.

\paragraph{Method.}
We route DWB's 2-bit tier through PolarQuant instead of scalar INT2:
\begin{enumerate}[leftmargin=*,topsep=2pt,itemsep=0pt]
  \item Apply WHT: $\tilde{\mathbf{k}} = W \mathbf{k}$ per head
        (head\_dim=64 $=$ 2$^6$, power-of-2 $\checkmark$)
  \item Apply uniform scalar quantization to $\tilde{\mathbf{k}}$
  \item Invert: $\hat{\mathbf{k}} = W^\top \tilde{\mathbf{k}}_q$
\end{enumerate}
All higher tiers (4/8/16-bit) use standard scalar quantization
unchanged.
The rotation is exactly invertible; no extra bits are consumed.

\paragraph{Results.}

\begin{table}[h]
\centering
\caption{\dwbtq\ results on SmolLM-360M, 100 samples.}
\label{tab:tq}
\footnotesize
\begin{tabular}{lcccc}
\toprule
Condition & HellaSwag & ARC-Ch. & avg\_bits \\
\midrule
FP16 & 41.0\% & 35.0\% & 16.0 \\
DWB-scalar & 40.0\% & 26.0\% & 5.05 / 7.72 \\
\textbf{\dwbtq} & \textbf{42.0\%} & \textbf{29.0\%} & 5.05 / 7.72 \\
Paper DWB & 41.2\% & --- & 5.05 \\
\bottomrule
\end{tabular}
\end{table}

\dwbtq\ recovers $+$2pp on HellaSwag and $+$3pp on ARC-Challenge
over DWB-scalar at \emph{identical compression}.
On HellaSwag, \dwbtq\ (42.0\%) exceeds the paper's DWB target (41.2\%)
and matches FP16 (42.6\%) within noise.
The gain is larger on ARC-Challenge ($+$3pp reasoning) than HellaSwag
($+$2pp commonsense), suggesting WHT rotation particularly benefits
reasoning-heavy token distributions.

ARC-Challenge bit distribution: $\{2: 37.4\%, 4: 17.3\%, 8: 12.2\%, 16: 33.2\%\}$—the
controller assigns more 16-bit on reasoning tasks, reducing the fraction
routed through WHT but yielding higher per-token gain.

%-----------------------------------------------------------------------
\section{Discussion}
\label{sec:discussion}
%-----------------------------------------------------------------------

\paragraph{On scale-dependent losslessness.}
Our mechanistic finding changes the interpretation of the H2 claim.
At 135M/360M, standard INT4 is already lossless—the paper's +7.6pp gain
is over a sub-standard baseline.
At 1.7B, standard INT4 genuinely degrades, and DWB's adaptive assignment
provides real value.
This scale-dependence follows from the hidden dimension and number of
attention heads, both of which affect error variance and cancellation.
Practitioners should evaluate their model's effective residual error
before assuming INT4 is a useful quantization target.

\paragraph{On the H3 implementation gap.}
Our best reproduction (v2, 37.0\%) is 4.2pp below the paper's 41.2\%.
The gap is primarily attributable to controller training quality.
The paper's compound loss with an explicit latency penalty is the key
mechanism that jointly optimizes accuracy and compression—something that
plain cross-entropy on importance quartiles cannot achieve.
The undisclosed training details (data distribution, exact loss weights,
number of training tokens) prevent full reproduction.
We expect the gap to close when the original code is released in June 2026.

\paragraph{On DWB-TurboQuant.}
Routing low-importance tokens through PolarQuant rather than naive scalar
INT2 provides a free accuracy improvement at no compression cost.
This suggests that the bit-width assignment and the \emph{quantization
method within each tier} are orthogonal dimensions that can be
independently optimized.
Future work could explore per-tier quantization methods—e.g., vector
quantization \cite{hooper2024kvquant} or product quantization for higher
bit tiers.

%-----------------------------------------------------------------------
\section{Conclusion}
\label{sec:conclusion}
%-----------------------------------------------------------------------

We re-implemented DWB \cite{haeri2026dwb} from paper equations and
conducted a systematic verification across SmolLM-135M, 360M, and 1.7B
on HellaSwag.
Our main finding is that symmetric INT4 KV quantization is
scale-dependent: lossless at $\leq$360M (effective residual 8.1\%) and
genuinely lossy at 1.7B (12.4\%), fully explained from first principles.
The paper's +7.6pp H2 claim is most valid at 1.7B where standard INT4
truly degrades.
For H3, we identify a dual-objective tension in controller training that
prevents simultaneously achieving the paper's accuracy \emph{and}
compression targets with our re-implementation; the undisclosed compound
loss is the key missing ingredient.
\dwbtq, our novel contribution, demonstrates that upgrading the 2-bit
scalar tier to WHT-based quantization recovers $+$2--3pp at no extra
cost—exceeding the paper's DWB target on HellaSwag.

\paragraph{Reproducibility statement.}
All code, data, and results are available at\\
\url{https://github.com/LonghornSilicon/dont-waste-bits}.
The latency experiment (H1) requires an RTX 4090 and is pending
hardware access; arithmetic verification confirms the 17.75\% reduction
is self-consistent with the paper's Table 3.

%-----------------------------------------------------------------------
\bibliographystyle{plain}
\bibliography{refs}
%-----------------------------------------------------------------------

\end{document}
